%! Unit 1: Algebraic Expressions and Factoring — Unit Packet Cover (KEY, 2pp)
%
% Page 1 is byte-identical to the student cover: both wrappers \input the same
% unit_cover/body.tex, so the two packets cannot drift.
%
% Page 2 is TEACHER-ONLY and is the reason this file exists. Per the teachernote
% convention, a unit test's answer rationale and Part D scoring must NOT sit at
% the foot of practice_test_key/actual_test_key: the practice test is bound into
% the STUDENT packet, so its rationale would ride along in front of students, and
% a note in a key also makes the key longer than its blank. shared/unit.mk merges
% this file into the KEY packet only. Keep the notes to ONE page — cover + notes
% is a single double-sided sheet.
\documentclass[10pt]{article}
\usepackage{atda-article}
\usepackage{atda-boxes}
\usepackage{ltablex}
\keepXColumns

\begin{document}

%! Unit 1: Algebraic Expressions and Factoring — the unit packet cover sheet.
%
% THIS FILE IS THE COVER. Both wrappers \input it — unit_cover/main.tex (student
% packet) and unit_cover_key/main.tex (key packet, which adds a page 2 of exam
% scoring notes) — so page 1 can never drift between the two packets. Edit the
% cover here, never in a wrapper.

% Full-bleed royal banner
\begin{tikzpicture}[remember picture, overlay]
  \fill[royal] (current page.north west)
    rectangle ([yshift=-0.9in]current page.north east);
\end{tikzpicture}
\vspace{-0.8in}

\begin{center}
  {\color{white}\textbf{\LARGE Algebra, Trigonometry, and Data Analysis}}\\[4pt]
  {\color{mistmid}\large\textbf{Unit 1: Algebraic Expressions and Factoring}}\\[6pt]
  {\color{white}\textbf{\large Unit Overview \quad\textbar\quad Eight Lessons \quad\textbar\quad Unit Test}}
\end{center}

\vspace{0.20in}
\namedateperiod
\vspace{0.14in}

\noindent
\small
Everything in this unit is one habit run in two directions. \textbf{Multiplying} takes two
factors and hides them inside a single expression; \textbf{factoring} takes that expression
apart and puts the factors back. Every lesson here is a different shape of that same move ---
and once an expression is factored, two new questions open up that you could not answer before:
\emph{when is it zero?} (Lesson 1.6) and \emph{what cancels?} (Lesson 1.7).

\vspace{0.12in}

% BOXGUARD: all three boxes below are breakable and all three fit on page 1 today
% with ~1.5in to spare, so these guards are inert as written. They are sized to
% their own boxes so that if the cover ever grows, a box moves whole to a second
% page instead of splitting and stranding a stub — the guard is the cheap
% insurance, and each was verified not to cost this sheet its single page.
\boxguard[14]
\begin{tocbox}
\small
\renewcommand{\arraystretch}{1.32}
\begin{tabularx}{\linewidth}{@{} c >{\raggedright\arraybackslash}p{5.5cm} l X @{}}
\rowcolor{mist}
\textbf{\#} & \textbf{Lesson} & \textbf{SOL} & \textbf{The one idea} \\
\midrule
1.0 & Unit Opener                    & ---                 & One model, two equivalent forms \\
1.1 & Expressions \& Exponent Rules  & \texttt{A2.EO.3a}   & An exponent counts factors, not terms \\
1.2 & Radicals \& Rational Exponents & \texttt{A2.EO.2a--c}& A radical is an exponent rule run backwards \\
1.3 & Factoring: GCF \& Grouping     & \texttt{A2.EO.3b}   & Pull out the \emph{largest} factor, not the first \\
1.4 & Factoring Trinomials           & \texttt{A2.EO.3b}   & Split the middle term, then group \\
1.5 & Factoring Special Forms        & \texttt{A2.EO.3b, d}& Patterns you already earned --- and can verify \\
1.6 & Solving by Factoring           & \texttt{A2.EI.2b, 6a--b} & Zero is the only product worth having \\
1.7 & Rational Expressions           & \texttt{A2.EO.1a--b}& Only \emph{factors} cancel \\
\end{tabularx}
\end{tocbox}

\vspace{0.12in}

\boxguard[16]
\begin{spiralbox}
\small
\begin{enumerate}[leftmargin=1.5em, itemsep=3pt, topsep=2pt]
  \item \textbf{An operation does not reach inside a sum.} $(2m+3)^2 \ne 4m^2+9$,
        $\sqrt{9+16} \ne 3+4$, and $\frac{x+5}{x+7} \ne \frac57$. Three lessons, one belief ---
        squaring, rooting, and cancelling all distribute over \emph{multiplication} only.
  \item \textbf{``Completely'' means largest, not first.} $6x(4x+6)$ is a true statement and a
        wrong answer. Check every factor you write down to see whether it factors again.
  \item \textbf{Factoring recovers information.} Multiplying erases the dimensions of a rectangle;
        factoring an area puts them back.
  \item \textbf{Zero is the only number that tells you anything.} $ab=12$ says nothing about $a$
        or $b$; $ab=0$ forces one of them to be zero. That is why every equation gets moved to
        standard form first.
  \item \textbf{A cancelled factor leaves a scar.} Excluded values are read from the
        \emph{original} denominator. Simplifying hides a restriction; it never removes one.
\end{enumerate}
\end{spiralbox}

\vspace{0.12in}

\boxguard[8]
\begin{remindbox}
\small
\textbf{The unit test has four parts and is worth 100 points:} Part A vocabulary matching
(8 pts), Part B multiple choice (12 pts), Part C short answer and computation (50 pts), and
Part D extended response (30 pts). A \textbf{practice test} with the same structure and
different numbers is bound at the back of this packet --- work it with no notes, then check it
against the key. Part D asks you to \emph{justify}, not just to compute: expect to verify an
identity by multiplying, and to say where two expressions agree and where they do not.
\end{remindbox}


\newpage

\begin{headlinebox}{royal}{\color{white}\bfseries Unit 1 --- Exam Scoring Notes (Teacher Copy)}\end{headlinebox}

\vspace{4pt}
\noindent\footnotesize
Both unit assessments are parallel forms: same blueprint, same ideas, different numbers and
reshuffled vocabulary letters. \textbf{P} = practice test (bound in the student packet),
\textbf{A} = actual test (distributed at test time only). 100 points each.

\vspace{4pt}

\boxguard[10]
\begin{teachernote}[Part A --- Vocabulary (8 pts, 1 each)]
\footnotesize
\textbf{P:} 1--C, 2--F, 3--J, 4--B, 5--G, 6--H, 7--D, 8--A. \qquad
\textbf{A:} 1--G, 2--D, 3--J, 4--H, 5--F, 6--A, 7--E, 8--B.\\[2pt]
Both forms carry \textbf{two distractor definitions} that match no term (P: E and I; A: C and I) ---
``factored form'' and ``exponent.'' They are there so a student who knows six terms cannot
back into the last two by elimination. Do not award part credit for a near miss on
\emph{prime polynomial}: ``cannot be factored'' without \emph{over the integers} is the whole
misconception behind $x^2-2$, and it is worth one sentence at review.
\end{teachernote}

\boxguard[10]
\begin{teachernote}[Part B --- Multiple Choice (12 pts, 2 each)]
\footnotesize
\textbf{P:} 1--C, 2--B, 3--C, 4--B, 5--B, 6--C. \qquad \textbf{A:} 1--C, 2--B, 3--C, 4--B, 5--B, 6--C.
(The keys match by design --- the two forms differ in numbers, not in option order, so a
student who memorized letters from the practice test gains nothing, because the stems changed.)\\[2pt]
\textbf{Items 1, 2 and 6 are one item asked three ways}, and that is the point: each distractor
(A) is the belief that squaring, rooting, or cancelling reaches inside a sum. A student who
misses all three has one problem, not three. \textbf{Item 3}'s wrong answers are all \emph{true}
statements that stop early --- score them zero, but name them ``incomplete,'' not ``wrong.''
\textbf{Item 5} needs only ``a square is never negative'' from the factor $x^2+k$, $k>0$; students
do not need complex arithmetic and should not be shown any.
\end{teachernote}

\boxguard[10]
\begin{teachernote}[Part C --- Short Answer \& Computation (50 pts, 5 each)]
\footnotesize
One item per lesson, in lesson order: C1 exponent rules (1.1), C2 radicals and rational
exponents (1.2), C3 grouping (1.3), C4 trinomial with $a\ne1$ (1.4), C5 special forms (1.5),
C6 solving (1.6), C7 simplifying (1.7), C8 dividing (1.7), C9 applied area, C10 factor
completely. \textbf{Suggested split: 3 pts method, 2 pts answer.}\\[2pt]
\textbf{Two deductions to apply consistently.} (i) On C7 and C8, a correct expression with
\emph{no excluded values stated} is 3 of 5 --- the restriction is half the standard
(\texttt{A2.EO.1b}), and on C8 the excluded value $x\ne-2$ (P) / $x\ne5$ (A) comes from a factor
that has already cancelled. (ii) On C10, a factorization that stops one step early
(e.g.\ $2(x^4-16)$) is 3 of 5. On C6, students may leave $x=-\tfrac32$ (P) / $x=-\tfrac23$ (A)
as a fraction --- do not require a decimal.
\end{teachernote}

\boxguard[10]
\begin{teachernote}[Part D --- Extended Response (30 pts, 15 each)]
\footnotesize
\textbf{D1 (\texttt{A2.EO.3d}, \texttt{A2.EO.1b}) --- 15 pts:} 5 pts the identity verified by
multiplying (all six products shown, middle terms cancelling); 5 pts the numerical test at both
values; 5 pts the sentence. \textbf{Full credit requires the student to say the algebra is
\emph{correct} and the claim is still \emph{wrong}} --- the two expressions agree everywhere
except at the excluded value, where the original is $\tfrac00$ and the simplified form returns a
number. A student who only writes ``you can't divide by zero'' without connecting it to the word
\emph{every} earns 3 of those 5.\\[2pt]
\textbf{D2 (\texttt{A2.EI.6a--b}, \texttt{A2.EI.2b}) --- 15 pts:} 4 pts factoring to
$-16(t-5)(t+1)$ (P) / $-16(t-7)(t+1)$ (A); 4 pts both zeros; 4 pts the two-sided interpretation;
3 pts number and type. \textbf{The two-sided interpretation is the assessed idea}: $t=-1$ is
rejected because it is before the launch, but in part (c) \emph{both} $t=0$ and $t=4$ (P) /
$t=6$ (A) are kept --- one is the launch, the other is the way down. A student who discards the
extra solution by reflex loses those 4 points; ``throw out the negative one'' is a judgement,
not a rule.
\end{teachernote}

\end{document}
